\section*{Notch 2}
\noindent\colorbox{orange}{
  \begin{minipage}{1.0\linewidth}
    \paragraph{Richiesta (a)}
    Utilizzando semplici argomenti di natura qualitativa (ovvero senza effettuare il calcolo
    della funzione di trasferimento, come richiesto al punto successivo) mostrare che il
    circuito svolga effettivamente la funzione di notch, discutendone la risposta ai limiti del
    dominio delle frequenze reali e verificando che i limiti asintotici del guadagno (sia per
    $f \to 0$ che per $f \to \infty$) coincidano con quanto osservato al punto 1b.
  \end{minipage}
}
\paragraph{Risposta (a)}
Qualitativamente il circuito si tratta di un amplificatore pilotato
dall'ingresso non invertente che viene preceduto da un nodo che
ridirige il flusso di corrente ad un blocco secondario da studiare
qualitativamente. Nei limiti di frequenze il sotto-circuito agisce
come se non esistesse. Il ramo in retroazione dell'ingresso invertente
agisce come un \textit{dimezzatore di tensione}, dunque se l'impedenza
effettiva del sotto-circuito equivale alla configurazione all'ingresso
invertente il guadagno differenziale sar\'a nullo e il filtro agira
come \textit{notch-filter}. Con uno studio quantitativo si potra calcolare
quanti notch sono presenti e dove sono disposti in funzione della frequenza.

\noindent\colorbox{orange}{
  \begin{minipage}{1.0\linewidth}
    \paragraph{Richiesta (b)}
    Derivare analiticamente la funzione di trasferimento (in trasformata di Laplace) A(s)
    nell'ipotesi di linearità (da giustificare a posteriori, vedi punto d) ed assumendo per
    semplicità che R1=R2=R', R3=R=20kW, R4=R5=R/2, C1=C=1nF e C2=100 C.
    Si suggerisce di procedere:
    \begin{itemize}
    \item[i]
      ottenendo un'espressione dell'impedenza di ingresso ZIN della rete inserita nel
      riquadro tratteggiato nello schema circuitale, ovvero quella vista dall'ingresso non
      invertente dell'operazionale U1 (\textit{suggerimento: si determini un'espressione della
        tensione all'ingresso non invertente dell'operazionale U1 in termini della corrente I1
        indicata in figura e considerando la relazione tra quella corrente e I2});
    \item[ii]
      ricavando una relazione tra VOUT, VS, R e ZIN sulla base dello schema equivalente
      riportato in calce ed infine calcolando il rapporto VOUT/VS in funzione di s (o, più
      convenientemente, della variabile complessa adimensionale $\xi=sRC$).
    \end{itemize}
  \end{minipage}
}
\paragraph{Risposta (b.i)}
Seguendo il suggerimento
\begin{equation*}
  I_1 \equiv sC(V_{+U1} - V_{+U2})
\end{equation*}
nell'ipotesi di op-amp ideale, +U2 non assorbe corrente, di conseguenza
\begin{equation*}
  I_1 = \frac{V_{+U2} - V_3}{R/2}
\end{equation*}
inoltre lo schema circuitale dell'op-amp U2 \'e quello di un inseguitore
di tensione, di conseguenza
\begin{equation*}
  V_{+U2} = V_{U2} \implies I_1 = \frac{V_{U2} - V_3}{R/2}
\end{equation*}
la corrente all'uscita dell'op-amp U2
\begin{equation*}
  I_2 \equiv 100sC(V_{U2} - V_3) \implies I_2 = 100sC\left(\frac{R}{2}I_1\right)
\end{equation*}
la KCL al nodo 3
\begin{equation*}
  \begin{gathered}
    I_1 + I_2 = I_3 \equiv \frac{V_3}{R/2} \\
    \implies V_3 = \frac{R}{2}(1 + 50sRC)I_1
  \end{gathered}
\end{equation*}
dalla seconda espressione di $I_1$
\begin{equation*}
  \begin{gathered}
    I_1 = \frac{V_{+U2} - \frac{R}{2}(1 + 50sRC)I_1}{R/2} \\
    \implies V_{+U2} = \frac{R}{2}(2 + 50sRC)I_1
  \end{gathered}
\end{equation*}
dalla prima espressione/definizione di $I_1$
\begin{equation*}
  \begin{gathered}
    I_1 = sC(V_{+U1} - R(1 + 25sRC)I_1) \\ \implies
    (1 + sRC + 25sR^2C^2)I_1= sCV_{+U1}
  \end{gathered}
\end{equation*}
l'impedenza d'ingresso $Z_{in}$ si ottiene dal rapporto fra la tensione
d'ingresso $V_{+U1}$ e la corrente d'ingresso $I_1$
\begin{equation}
  \label{eq:notch2-zin}
  \colorbox{blu}{\boxed{Z_{in} \equiv \frac{V_{+U1}}{I_1} = \frac{1}{sC}(1 + sRC + 25sR^2C^2)}}
\end{equation}

\paragraph{Risposta (b.ii)}
La tensione all'ingresso invertente $V_-$ \'e met\'a della tensione sorgente $V_s$
(partitore di tensione con resistenze uguali dimezzano la tensione)
\begin{equation*}
  V_+ = \frac12V_s + \frac12V_{OUT}
\end{equation*}
la tensione all'ingresso \underline{non}-invertente risulta sempre come risultato
di un partiore composto da $R$ e $Z_{in}$
\begin{equation*}
  V_- = \frac{Z_{in}}{R + Z_{in}}V_s
\end{equation*}
nell'ipotesi di op-amp ideale che opera in regime di linearit\'a si applica il principio
di corto-circuito virtuale
\begin{equation*}
  V_+ = V_- \implies \frac12V_s + \frac12V_{OUT} = \frac{Z_{in}}{R + Z_{in}}V_s
\end{equation*}
la forma generale della funzione di trasferimento parametrizzata da $Z_{in}$
\begin{equation*}
  H(s) = \frac{Z_{in} - R}{Z_{in} + R}
\end{equation*}
espandendo l'espressione di $Z_{in}$ si ottiene
\begin{equation}
  \label{eq:notch2-tfunc}
  \colorbox{blu}{\boxed{H(s) = \frac{1 + 25\xi^2}{1 + 2\xi + 25\xi^2} \qquad \xi \equiv sRC}}
\end{equation}

\noindent\colorbox{orange}{
  \begin{minipage}{1.0\linewidth}
    \paragraph{Richiesta (c)}
    Verificare che l'espressione ottenuta per la funzione di trasferimento soddisfi la
    condizione necessaria a garantire la stabilità del filtro in zona lineare.
  \end{minipage}
}
\paragraph{Risposta (c)}
La stabilit\'a del filtro si verifica studiando i poli della funzione di
trasferimento complessa $H(s)$ (ovvero gli zeri del denominatore)
\begin{equation*}
  \begin{gathered}
    1 + 2\xi + 25\xi^2 = 0 \implies \xi = \frac{-2 \pm i\sqrt{96}}{50} \\ \implies
    \colorbox{blu}{\boxed{\mathcal{R}\{\xi\} = -\frac{1}{25} < 0}}
  \end{gathered}
\end{equation*}
la parte reale dei poli \'e negativa, per il terzo criterio di stabilit\'a
questo implica che il filtro \'e stabile e lineare. \\

\noindent\colorbox{orange}{
  \begin{minipage}{1.0\linewidth}
    \paragraph{Richiesta (d)}
    Dedurre i valori attesi (senza incertezze) per i parametri del filtro (frequenza centrale,
    larghezza di banda e fattore di merito) e confrontarli con le misure di cui al punto 1b (si
    consiglia di esprimere la funzione di trasferimento sul dominio delle frequenze reali in
    termini della grandezza adimensionale $u = \omega RC$).
  \end{minipage}
}
\paragraph{Risposta (d)}
Seguendo il consiglio della richiesta per la scelta di parametri, le parti reali
e immaginarie del numeratore e denominatore della funzione di trasferimento
complessa in forma canonica
\begin{equation*}
  \begin{cases}
    \mathcal{R}_{\omega}^N = 1 - 25u^2 \\
    \mathcal{I}_{\omega}^N = 0 \\
    \mathcal{R}_{\omega}^D = 1 - 25u^2 \\
    \mathcal{I}_{\omega}^D = 2u
  \end{cases}
\end{equation*}
il guadagno \'e ottenuto dall'espressione \ref{eq:da-tfunc-a-gain}
\begin{equation*}
  G(u) = \frac{1 - 25u^2}{\sqrt{(1 - 25u^2)^2 + (2u)^2}}
\end{equation*}
nel contesto attuale, ovvero nello studio di un filtro notch, la frequenza
centrale \'e quella frequenza per cui il guadagno si annulla (il \textit{notch})
\begin{equation*}
  1 - 25u^2 = 0
\end{equation*}
sostituendo $u \to \omega RC$ e $\omega \to 2\pi f$ si ottiene la frequenza centrale
\begin{equation}
  \label{eq:notch2-frequenza-centrale}
  \colorbox{blu}{\boxed{f = \frac{1}{10\pi RC}}}
\end{equation}
per studiare la larghezza di banda servono le frequenze di taglio e di conseguenza
anche il guadagno massimo, ri-arrangiando l'espressione del guadagno in termini di $u$
(nell'ipotesi che il numeratore non sia nullo quindi diverso dalla frequenza centrale)
si ottiene
\begin{equation*}
  G(u) = \frac{1}{\sqrt{1 + \left(\frac{2u}{1 - 25u^2}\right)^2}}
\end{equation*}
la funzione di trasferimento \'e il reciproco della radice di un espressione
quadratica con minimo uguale ad $1$ per valori reali, dunque il guadagno massimo
\begin{equation*}
  \max_{u \in [0, +\infty)} G(u) = 1
\end{equation*}
ora le frequenze di taglio si determinano come quelle frequenze per cui il guadagno
quadro \'e uguale alla met\'a del suo massimo quadro
\begin{equation*}
  \begin{gathered}
    \frac{1 - 25u^2}{\sqrt{(1 - 25u^2)^2 + (2u)^2}} = \frac{1}{\sqrt2} \\ \implies
    (1 - 25u^2)^2 = (2u)^2 \implies
    \begin{cases}
      1 - 25u_+^2 = +2u_+ \\
      1 - 25u_-^2 = -2u_-
    \end{cases} \\ \implies
    u_+ = \frac{-1 \pm \sqrt{26}}{25} \quad \wedge \quad u_- = \frac{1 \pm \sqrt{26}}{25}
  \end{gathered}
\end{equation*}
una frequenza ammissibile deve essere positiva, di conseguenza
\begin{equation*}
  u_{L,H} = \frac{\sqrt{26} \pm 1}{25}
\end{equation*}
attraverso la sostituzione $u \to \omega RC$ e $\omega \to 2\pi f$ si ottengono
le frequenze di taglio
\begin{equation*}
  f_{L,H} = \frac{\sqrt{26} \pm 1}{50\pi RC}
\end{equation*}
dunque la larghezza di banda
\begin{equation}
  \label{eq:notch2-bw}
  \colorbox{blu}{\boxed{\Delta f = \frac{1}{25\pi RC}}}
\end{equation}
infine il fattore di merito (Q-factor) si ottiene dal rapporto fra \ref{eq:notch2-frequenza-centrale}
e \ref{eq:notch2-bw}
\begin{equation}
  \label{eq:notch2-q-factor}
  \colorbox{blu}{\boxed{Q = \frac{25\pi RC}{10\pi RC} = \frac52}}
\end{equation}
questo valore \'e indipendente dai componenti del circuito. \\

\noindent\colorbox{orange}{
  \begin{minipage}{1.0\linewidth}
    \paragraph{Richiesta (e)}
    Discutere gli andamenti osservati nei plot di Bode di cui al punto 1a e
    giustificarli alla luce delle proprietà di simmetria della funzione di trasferimento.
  \end{minipage}
} \smallskip \\
\paragraph{Risposta (e)}

\noindent\colorbox{orange}{
  \begin{minipage}{1.0\linewidth}
    \paragraph{Richiesta (f)}
    Inviata in VS un'onda quadra di frequenza pari ad f0 ed ampiezza 500mV,
    osservare all'oscilloscopio ingresso (C1), uscita (C2) e la loro differenza (a tale scopo
    occorre aggiungere un ulteriore canale mediante il menù a tendina ``Add Channel →
    Math → Simple (software)'' associando la differenza dei due canali) e discutere
    qualitativamente e quantitativamente le caratteristiche di quest'ultima.
  \end{minipage}
} \smallskip \\
\paragraph{Risposta (f)}
